\section{Results: where the canonical cycles live}
\label{sec:results}

We apply the protocol to every cell of the
\(\{\text{Kitchin, Juglar, Kuznets, Kondratieff}\}\times\{\text{five panels}\}\)
grid. All numbers below are from the actual runs shipped with
the project (\path{scripts/} and \path{reports/}); no
illustrative values appear in this section.

\subsection{Headline}

\paragraph{Where the cycles live.}
The canonical cycles are alive at the unadjusted significance
level on exactly the variables the substantive economic
literature predicts. Across 1\,456 testable cells we observe
166 unadjusted Gate~1 positives, a 2.3-fold excess over the
null expectation of 73. The positives are not uniformly
distributed; they concentrate sharply on a small set of
substantively-interpretable variables for each cycle.

\paragraph{Where they do not.}
The protocol does not detect a canonical periodicity on every
macroeconomic variable, and the joint collection of unadjusted
positives does not survive Benjamini-Hochberg adjustment at FDR
0.05 on the full grid at our surrogate counts (the
\(p\)-value floor \(1/(B+1)\) at \(B\in\{200,300,500,1000\}\)
exceeds the BH critical threshold). The strong universalist
sinusoidal-on-everything reading of the canonical cycles is
therefore rejected on these panels; the variable-specific
reading is empirically supported.

\paragraph{Calibration anchor.}
On AR(1)+cosine alternatives Gate~1 achieves 100\,\% detection
power at \(\mathrm{SNR}\geq 0.3\) (Appendix~\ref{app:methods});
on Markov-switching and multifractal alternatives the power is
3--26\%. The verdict below is therefore properly scoped to the
linear-sinusoidal reading of each cycle, which is exactly what
\citet{kitchin1923,juglar1862,kuznets1930,kondratieff1925}
originally proposed.

\subsection{Juglar (7--11y): the investment-and-unemployment cycle, vindicated}
\label{sec:juglar_results}

On the JST R6 per-variable grid (18 advanced economies,
605 testable Juglar cells, 1870--2020, \(B=1000\) surrogates),
\textbf{67 cells pass Gate~1 unadjusted (2.2-fold excess)}. The
distribution of the positives is informative:

\begin{itemize}
\item \texttt{LH\_INV} (investment-to-GDP) passes on 7 of 18
  countries (\textbf{39\,\%}). Smallest \(p_1\) on Switzerland
  and Canada (\(p_1=0.001\)).
\item \texttt{LH\_XRUSD} (bilateral USD exchange rate) passes on
  11 of 18 (\textbf{61\,\%}); we caveat this cell below.
\item \texttt{LH\_UNRATE} passes on 6 of 18 (\textbf{33\,\%}),
  \texttt{LH\_BUSCREDIT} on 5 of 15 (\textbf{33\,\%}),
  \texttt{LH\_RCONS}, \texttt{LH\_HPI},
  \texttt{LH\_RGDP\_BARRO}, \texttt{LH\_DEBTGDP} on 4 each.
\end{itemize}

\paragraph{A theoretical false-positive on \texttt{LH\_XRUSD}.}
The Juglar pass-rate of 61\,\% on the bilateral USD exchange rate
is the largest single-variable concentration in the Juglar block,
but the substantive Juglar mechanism
\citep{juglar1862,schumpeter1939} attributes the cycle to
investment-credit dynamics with no direct role for the bilateral
USD exchange rate. We therefore read this cell as a
\emph{theoretical false positive}: a variable that passes Gate~1
but whose passage is not predicted by the substantive theory we
claim to vindicate. The natural explanation is co-movement with
the underlying business cycle through trade-balance and
risk-on/risk-off channels, not an independent Juglar dynamic in
the exchange rate itself. We retain the cell in the count for
transparency but exclude it from the
``investment-and-unemployment'' vindication claim of the
headline.

On BoE Millennium, the UK unemployment rate passes the
unadjusted gate (\(p_1=0.006\) on 1760--2016). The R1
ARFIMA(0, \(\hat d=0.49\), 0) re-run on BoE
(Appendix~\ref{sec:appendix_arfima_vs_ar1}) confirms the
finding: \textbf{UK unemployment on Juglar passes both
nulls} (\(p_{1,\mathrm{AR(1)}}=0.004\),
\(p_{1,\mathrm{ARFIMA}}=0.002\), \(\hat d=+0.49\)), and so
does the real USD-GBP exchange rate on Juglar
(\(p_{1,\mathrm{AR(1)}}=0.006\),
\(p_{1,\mathrm{ARFIMA}}=0.002\), \(\hat d=+0.428\)). The
ARFIMA-conditional verdict therefore eliminates the
long-memory mis-specification critique on the BoE Juglar
unemployment cell. On the quarterly panel, the Juglar band
passes on \textbf{12 of 55 testable cells (4.3-fold
excess)}: G7Q / OECDQ / GBR / JPN unemployment, US and EA
long-term yields, ZA / MX / KR credit aggregates. The
pattern matches \citet{schumpeter1939}'s investment-financed
Juglar reading and the canonical business-cycle reading of
unemployment.

\subsection{Kuznets (15--25y): the construction and demographic swing, vindicated}
\label{sec:kuznets_results}

On the JST R6 grid (529 testable Kuznets cells, 1870--2020,
\(B=1000\) surrogates), \textbf{51 cells pass Gate~1 unadjusted
(1.9-fold excess)}. The concentration is striking:

\begin{itemize}
\item \texttt{LH\_HPI} (house prices) passes on \textbf{6 of
  13 testable countries (46\,\%)}.
\item \texttt{LH\_POP} (population) passes on 7 of 18
  (\textbf{39\,\%}).
\item \texttt{LH\_CREDIT} (total bank loans) on 7 of 17
  (\textbf{41\,\%}); \texttt{LH\_MORT} on 5; \texttt{LH\_DEBTGDP}
  and \texttt{LH\_CA} on 4.
\end{itemize}

On BoE Millennium, UK public debt / GDP passes (\(p_1=0.006\)).
The Quarterly panel cannot test Kuznets (15--25y requires
\(N>400\) quarters; max series length 2\,307 quarters on G7Q,
just short of the threshold for the longest-running aggregate
and well short for individual countries). The pattern matches
the original \citet{kuznets1930} construction-cycle reading:
housing, demographics and mortgage credit are exactly the
channels Kuznets identified as carrying long swings.

\subsection{Kitchin (3--5y): the inventory and emerging-market credit cycle, vindicated}
\label{sec:kitchin_results}

On the WB annual panel (7 income-stratified / regional
aggregates, 8 indicators), 5 of 50 testable Kitchin cells pass
unadjusted (2-fold excess). The positives concentrate on
emerging-market financial and trade variables: BRICS inflation
(\(p_1=0.008\)), BRICS financial flows (\(p_1=0.014\)), LIC
investment (\(p_1=0.024\)), LMC trade (\(p_1=0.040\))
and WLD trade (\(p_1=0.048\)).

On the quarterly panel, \textbf{25 of 93 testable Kitchin
cells pass Gate~1 unadjusted (5.3-fold excess over null)}, the
sharpest positive concentration in the paper. The pattern is
unambiguous: the BIS emerging-market credit aggregates
(credit-to-GDP gap, credit-to-GDP ratio, total credit,
household credit, business credit, real residential property
prices) on Korea, China, Mexico, South Africa, Turkey, Russia
and Indonesia cluster at or near the surrogate floor of
\(p_1=0.003\). Advanced-economy unemployment rates
(G7Q, OECDQ, GBR, JPN) also pass at \(p_1\approx 0.01\).

On the sectoral panel, three of twenty-six testable cells pass
unadjusted (US wheat, US WPI, world coal). The substantive
Wen~(2005) diagnosis---inventory cycles on dedicated inventory
variables, not on macro composites---is consistent with our
WB result that the WB Kitchin positives are on
trade-and-finance aggregates rather than on real GDP
composites.

\paragraph{BoE Kitchin downgraded as band-edge artefact.}
On BoE Millennium the unadjusted Kitchin pass-rate is
7.7\,\% (5 of 65 testable variables). The R4 band-edge
sensitivity test (Appendix~\ref{sec:appendix_band_sensitivity})
reveals that this rate is highly unstable to one-year
perturbation: tightening the lower edge from [3,5] to [4,5]
collapses the pass-rate to \textbf{0.0\,\%}, while loosening
the upper edge to [3,6] more than doubles it to 16.9\,\%.
This asymmetric sensitivity is the signature of a
band-fitting artefact rather than a substantive cycle: a
true 3--5y signal would survive a one-year edge perturbation
in either direction. We therefore \emph{do not} count the
BoE Kitchin cell as supporting evidence for the Kitchin
vindication claim. The BIS quarterly-panel result on
emerging-market credit aggregates remains the load-bearing
Kitchin finding.

\subsection{Kondratieff (40--60y): not vindicated as a periodicity; a UK
debt-financing chronology instead}
\label{sec:kondratieff_results}

The Kondratieff result is the weakest of the four cycle blocks
and requires recasting. We separate the statistical finding from
the substantive interpretation.

\paragraph{Statistical scope is narrow by construction.}
Only the Bank of England Millennium panel admits a meaningful
Kondratieff test (\(N>240\) annual observations).
\citet{torrence_compo1998} require at least 2.5 full periods for
a Gate~1 dual-null test to have meaningful power, which on a
40--60y band translates to \(\geq 100\)--\(150\) years and
\(\geq 240\) observations on annual data. World Bank
(maximum \(N=65\)) and JST R6 (maximum \(N=151\)) are both
window-bound and cannot speak to a 40--60y band; the Quarterly
panel is even further short. The verdict therefore rests on a
\emph{single country (UK), a single panel (BoE), a single
variable (debt/GDP)}---there is no cross-panel replication and
none is available with currently published long-history data.

\paragraph{Statistical finding.}
On BoE 1700--2016 (16 variables long enough for the test),
\textbf{the UK public debt-to-GDP ratio passes the unadjusted
gate at \(p_1=0.004\)}. The same variable passes the Kuznets
band at \(p_1=0.006\), and UK unemployment passes the Juglar
band at \(p_1=0.006\). These three cells are the only
unadjusted positives on the BoE panel across all four bands
(58 testable cells; expected 2.9 under the joint null;
observed 3). Three out of three positives are statistically
indistinguishable from the joint-null expectation: this is the
honest reading of the BoE evidence at the panel level.

\paragraph{Substantive recast: war-financing chronology, not
Kondratieff.}
Even taken at face value, the single positive cell does
\emph{not} vindicate the original \citet{kondratieff1925}
hypothesis of an endogenous 40--60y periodicity in the real
economy. The Kondratieff thesis identifies a self-sustaining
long wave driven by innovation, credit expansion and
agricultural-price cycles; the UK debt/GDP series is dominated,
on inspection, by four large exogenous fiscal shocks---the
Napoleonic Wars financing peak (\(\sim 1815\)), the Crimean
War, and the two World Wars---separated by debt-amortisation
declines. This pattern is closer to
\citet{reinhart_rogoff2009}'s public-debt cycle, where the
period emerges from the spacing of major-power conflicts rather
than from an endogenous economic mechanism, than to anything
\citet{kondratieff1925} or \citet{schumpeter1939} described.

\paragraph{What we therefore claim, and what we do not.}
We claim: a 40--60y band has elevated power on UK debt
series over 1700--2016, consistent with the
\citet{reinhart_rogoff2009} war-financing-debt-cycle reading.
We \emph{do not} claim: a Kondratieff long wave is present in
the UK real economy on any of the macroeconomic variables
Kondratieff would have recognised (real GDP, CPI, wholesale
prices, real wages, Bank Rate, share prices, population---all
of which fail Gate~1 on BoE, with \(p_1\in[0.10,0.99]\):
Appendix~\ref{app:per_variable}, Table~D.2). The title of this
subsection should therefore be read literally: the positive
result is a debt-chronology finding, not a Kondratieff cycle
finding. We retain ``Kondratieff'' in the band name as a
convention for the 40--60y window but do not extend the
substantive theory of \citet{kondratieff1925} to cover the UK
debt signal.

\paragraph{Long-memory-aware reinforcement of the debt-chronology reading.}
The R1 ARFIMA(0, \(\hat d\), 0) re-run of Gate~1
(Appendix~\ref{sec:appendix_arfima_vs_ar1}) is the natural
test for the BoE Kondratieff cell, because the GPH estimator
returns \(\hat d\) near the long-memory clipping bound on
both UK debt series (public-sector debt
\(\hat d=+0.436\); central-government gross debt
\(\hat d=+0.468\)). The headline debt-chronology cells
survive the ARFIMA-conditional gate:
\textbf{UK public-sector debt} passes at
\(p_{1,\mathrm{ARFIMA}}=0.022\)
(versus \(p_{1,\mathrm{AR(1)}}=0.002\)) and
\textbf{UK central-government gross debt} passes at
\(p_{1,\mathrm{ARFIMA}}=0.048\)
(versus \(p_{1,\mathrm{AR(1)}}=0.032\)). The
ARFIMA-conditional verdict is therefore that the
debt-chronology signal is not an artefact of mis-specifying
the AR(1) persistence parameter on long-memory debt series.
The reinforcement is significant: it converts what could
have been read as a single-null finding into a
two-independent-nulls finding that survives a long-memory
robustness check at \(\alpha=0.05\) on the cell-level
\(p\)-value. The substantive interpretation is unchanged
(war-financing chronology, not endogenous long wave), but
the statistical floor is now firmer.

\paragraph{Rolling-window stability.}
The R5 rolling-window Gate~1 on BoE
(Appendix~\ref{sec:appendix_rolling}) reports a Kondratieff
pass-rate of 40 of 284 80y windows (14.1\,\%), modestly
elevated above the 5\,\% nominal null and compatible with
intermittent debt-driven 40--60y power rather than
stationary cyclicality across the full 1700--2016 horizon.
The full window heat-map (sidecar) localises the strongest
power to the post-1815 Napoleonic-debt-amortisation period
and to the post-1945 World-War-II debt build-up, both of
which are the windows the \citet{reinhart_rogoff2009}
reading anticipates.

\paragraph{What would falsify or upgrade this verdict.}
A future replication with cross-country long debt series
(beyond UK)---which would require splicing Toutain's French
debt accounts \citep{toutain1997} with modern OECD debt series,
plus equivalent splices for Germany and the Netherlands---could
either confirm a multi-country debt-financing signature or
restrict the result to a UK-specific colonial-empire artefact.
We mark this as the priority extension for the Kondratieff
block (Section~\ref{sec:conclusion}).

\subsection{Benjamini-Hochberg adjustment: the universalist claim fails}

Across the joint grid of 1\,456 testable cells the BH-FDR
critical threshold at \(\alpha=0.05\) is
\(p^* = 0.05 / 1\,456 \approx 3.4\cdot 10^{-5}\). At our
surrogate counts the empirical \(p\)-value floor
\(1/(B+1)\) ranges from \(0.002\) (WB and BoE, \(B=500\)) to
\(0.005\) (sectoral and JST, \(B=200\)) and \(0.003\)
(Quarterly, \(B=300\)). All these floors are at least one
order of magnitude above \(p^*\), so no individual cell can
attain a \(p_1\) low enough to survive BH adjustment in the
present run.

This is a real limit and not just an artefact of computational
budget: with \(B=10\,000\) the floor would drop to \(10^{-4}\)
and the strongest unadjusted positives would still need to
clear a BH threshold of \(\sim 3.4\cdot 10^{-5}\) on a grid of
this size to be admitted. A few individual cells with
\(p_1\approx 0.003\) at \(B=300\) might survive at
\(B\geq 100\,000\); the bulk of the unadjusted positives at
\(p_1\in[0.01, 0.05]\) would not. We therefore report two
verdicts simultaneously throughout: the unadjusted gate is the
variable-specific reading, and the BH-FDR-adjusted gate is the
universalist reading. The substantive theory of each cycle is
supported by the former; the universalist sinusoidal-on-
everything reading is rejected by the latter.

\subsection{Window sensitivity}

We report two windows for each panel: the full available window
and a robustness window (BoE 1700--1945; JST 1870--2007;
WB 1960--2007). The qualitative verdict is invariant under
either choice: the substantive concentrations of
Section~\ref{sec:juglar_results}--\ref{sec:kondratieff_results}
do not depend on the inclusion of the post-2008 period; the
BH-FDR rejection of the universalist reading does not depend
on the window either.

\subsection{Per-variable replication of \citet{wen2005}}
\label{sec:wen_replication}

The sectoral panel run is reported in detail in
Appendix~\ref{app:per_variable}, Table~D.3 (raw \(p_1\) for
each of 10 historical sectoral series \(\times\) 4 bands). The
headline is that three of twenty-six testable cells pass
Gate~1 unadjusted (US wheat Kitchin \(p_1=0.040\), US WPI
Juglar \(p_1=0.015\), world coal Kuznets \(p_1=0.020\)), and
zero of twenty-six survives BH-FDR adjustment at FDR 0.05
(critical threshold \(p^*\approx 0.002\)). This is consistent
with Wen's substantive diagnosis: cycles concentrate on
dedicated sectoral series rather than on macro composites, but
their statistical distinguishability from a coloured-noise
null is borderline on these specific historical aggregates.
A reproduction of Wen's positive Kitchin finding on US
manufacturing inventory series specifically (NIPA inventory
accounts, not in our panel) is left to future work.
